\documentclass[12pt]{article}
\usepackage[utf8]{inputenc}
\usepackage[T1]{fontenc}
\usepackage{hyperref}
\title{Ateso Language Resources for K-12 Curriculum Development}
\author{Compiled by Claude Sonnet 5 (claude-sonnet-5) for Andr\'as Kornai}
\date{}
\begin{document}
\maketitle

\section{Introduction}
Ateso (ISO 639-3 teo) is a Nilotic language of Uganda and Kenya. This report is based on a single web search, which returned mostly bibliographic leads and no freely available primer or textbook. Wikipedia's figure of 2,778,000 native speakers (2014 and 2019 censuses) is higher than the roughly 1,900,000 in the task description.

\section{Primary Educational Resources}
No online school textbook or primer was found. Print primers are listed in bibliographies.
\begin{itemize}
\item Hilders and Lawrance, \emph{An Introduction to the Ateso Language}, Eagle Press, Kampala. WALS dates it to 1956 and SOAS records to 1957. I could not confirm whether it is in print or online. \url{https://wals.info/refdb/record/Hilders-and-Lawrance-1956}
\end{itemize}

\section{Linguistic and Reference Materials}
\begin{itemize}
\item Bibliography of Teso lexicons and grammars in the CNRS REFLEX database. It lists the Ongodia and Ejiet Ateso--English dictionary (2008), Kiggen's Ateso--English dictionary (1953), and a study of Ateso derivational verbal morphology by Otaala (1981). \url{https://reflex.cnrs.fr/Lexiques/webball/biblio.php?LGNOM=Teso}
\item SOAS Library record that appears to cover the Hilders and Lawrance English--Ateso and Ateso--English vocabulary (1958). It also holds unpublished typescripts by A. W. Henriksen (1957) on the Teso verb and the preposition in Ateso. The search result title is only a reference number, so this is a best reading. \url{https://library.soas.ac.uk/Collection/MS_380573}
\item Wikipedia article on Teso language. It mentions the 1947 orthography committee and a later orthography report (1964). \url{https://en.wikipedia.org/wiki/Teso_language}
\end{itemize}

\section{Religious and Community Resources}
None were found in this search.

\section{Digital and Audio Resources}
None were found in this search.

\section{Curriculum Development Context}
The search found only older print materials and bibliographic records. A curriculum project would need to check library catalogues (WorldCat, SOAS, Makerere University library) and ask Ugandan education authorities and publishers about current Ateso school materials.
\end{document}